\section{Pattern と matcher の規則}

\subsection{ユーザーパターンの synthesis}

$\DDPattern$ は pattern dual $\eta=\kappa\triangleright\tau$ と binding context を
同時に合成する．$\Delta$ は左から右へ延長し，同じ branch 内の重複 binder を拒否する．
\begin{mathpar}
\inferrule{
  x\notin\dom(\Delta)
}{
  \DDPattern_{\Sighat}(q,S,\Gamma,\Phi,\Delta,\$x,
    \chi_{q_\kappa}\triangleright\alpha_{q_\tau},
    \Delta,x:\alpha_{q_\tau},q+\langle1,1\rangle,S)
}\rn{DD-PAT-VAR}

\inferrule{ }{
  \DDPattern_{\Sighat}(q,S,\Gamma,\Phi,\Delta,\_,
    \chi_{q_\kappa}\triangleright\alpha_{q_\tau},\Delta,
    q+\langle1,1\rangle,S)
}\rn{DD-PAT-WILD}

\inferrule{
  q;S;\mono\Delta,\Gamma\vdash e\Rightarrow\tau\dashv q_1;S_1
}{
  \DDPattern_{\Sighat}(q,S,\Gamma,\Phi,\Delta,\#e,
    \chi_{(q_1)_\kappa}\triangleright\tau,\Delta,
    q_1+\langle1,0\rangle,S_1)
}\rn{DD-PAT-VALUE}

\inferrule{
  \Phi(x)=\kappa\triangleright\tau
}{
  \DDPattern_{\Sighat}(q,S,\Gamma,\Phi,\Delta,\widetilde x,
    \kappa\triangleright\tau,\Delta,q,S)
}\rn{DD-PAT-EMBED}
\end{mathpar}

tuple と constructor は子 pattern を先に合成する．constructor は instantiated field target と
子 target を align し，観測可能位置から result capability を決め，最後に
$\capcompat$ を検査する．
\begin{mathpar}
\inferrule{
  \DDPatterns_{\Sighat}(q,S,\Gamma,\Phi,\Delta,
    \overline p,\overline\eta,\Delta',q',S')
}{
  \DDPattern_{\Sighat}(q,S,\Gamma,\Phi,\Delta,(\overline p),
    (\overline{\eta.\kappa})\triangleright
    (\overline{\eta.\tau}),\Delta',q',S')
}\rn{DD-PAT-TUPLE}

\inferrule{
  \Inst_q(\Sighat_P(c))=(\overline\rho\Rightarrow_P\rho,q_0) \\
  \DDPatterns_{\Sighat}(q_0,S,\Gamma,\Phi,\Delta,
    \overline p,\overline\eta,\Delta',q_1,S_1) \\
  \mathsf{DDAlignTargetList}(S_1,\overline\eta,\overline\rho,S_2) \\
  \mathsf{DDPatternCtorCap}_{\Sighat}
    (c,q_1,S_2,\overline{\eta.\kappa},\kappa,q_2,S_3) \\
  \capcompat_{\Sighat}(c;S_3\overline{\eta.\kappa},S_3\kappa)
}{
  \DDPattern_{\Sighat}(q,S,\Gamma,\Phi,\Delta,c\,\overline p,
    \kappa\triangleright\rho,\Delta',q_2,S_3)
}\rn{DD-PAT-CON}
\end{mathpar}

$p_1\mathbin\&p_2$ は左の bindings を右へ渡す．$p_1\mathbin|p_2$ は同じ入力
$\Delta$ から両 alternative を調べ，dual と binder 名・binder type を最後に align する．
\begin{mathpar}
\inferrule{
  \DDPattern(q,S,\Gamma,\Phi,\Delta,p_1,\eta_1,
    \Delta_1,q_1,S_1) \\
  \DDPattern(q_1,S_1,\Gamma,\Phi,\Delta_1,p_2,\eta_2,
    \Delta_2,q_2,S_2) \\
  \mathsf{DDAlignDual}(S_2,\eta_1,\eta_2,S')
}{
  \DDPattern(q,S,\Gamma,\Phi,\Delta,p_1\mathbin{\&}p_2,
    \eta_1,\Delta_2,q_2,S')
}\rn{DD-PAT-AND}

\inferrule{
  \DDPattern(q,S,\Gamma,\Phi,\Delta,p_1,\eta_1,
    \Delta_1,q_1,S_1) \\
  \DDPattern(q_1,S_1,\Gamma,\Phi,\Delta,p_2,\eta_2,
    \Delta_2,q_2,S_2) \\
  \mathsf{DDAlignDual}(S_2,\eta_1,\eta_2,S_3) \\
  \mathsf{DDAlignBindings}(S_3,\Delta_1,\Delta_2,S')
}{
  \DDPattern(q,S,\Gamma,\Phi,\Delta,p_1\mathbin{|}p_2,
    \eta_1,\Delta_1,q_2,S')
}\rn{DD-PAT-OR}
\end{mathpar}

pattern-function scheme は freeze 済み signature から instantiate する．declaration の検査は
source expression の $\mathsf{DDTyping}$ の外側であり，本稿に declaration typing rule は置かない．
\begin{mathpar}
\inferrule{
  \InstDual_q(\Sighat_F(f))=(\overline\eta\Rightarrow\eta,q_0) \\
  \DDPatterns(q_0,S,\Gamma,\Phi,\Delta,\overline p,
    \overline\eta',\Delta',q_1,S_1) \\
  \mathsf{DDAlignDualList}(S_1,\overline\eta',\overline\eta,S')
}{
  \DDPattern(q,S,\Gamma,\Phi,\Delta,f\,\overline p,
    \eta,\Delta',q_1,S')
}\rn{DD-PAT-APP}
\end{mathpar}

\subsection{MatchAll の型付け}

$\mathtt{matchAll}$ は target を合成し，pattern target と align した後，matcher expression を
slot type に対して check する．第四 premise は，\textsc{DD-MATCHALL} 自身が直接導入する唯一の
slot demand である．ただし，これは導出木全体に demand が一つしかないという意味ではない．
例えば $e_m$ が matcher literal なら，その synthesis 中に各 clause の next matcher が
それぞれ別の slot に対して check される．各 checking cut が独立に，高々一度の coercion を行う．
\begin{mathpar}
\inferrule{
  q;S;\Gamma\vdash e_t\Rightarrow\tau_t\dashv q_1;S_1 \\
  \DDPattern(q_1,S_1,\Gamma,\varnothing,\varnothing,p,
    \kappa_p\triangleright\tau_p,\Delta,q_2,S_2) \\
  \DDAlignTypes(S_2,\tau_p,\tau_t,S_3) \\
  q_2;S_3;\Gamma\vdash e_m\Leftarrow
    \MatcherSlot{\kappa_p}{\tau_t}\dashv q_3;S_4 \\
  q_3;S_4;\mono\Delta,\Gamma\vdash e\Rightarrow
    \tau_r\dashv q';S'
}{
  q;S;\Gamma\vdash
  \mathtt{matchAll}\ e_t\ \mathtt{as}\ e_m\
  \mathtt{with}\ p\to e\Rightarrow
  \ListTy{\tau_r}\dashv q';S'
}\rn{DD-MATCHALL}
\end{mathpar}

\subsection{Primitive pattern と data pattern の検査}

$\DDPPat$ は一つの shared matcher target に対して primitive pattern を check し，hole dual 列と
bindings を返す．hole だけが fresh consumer capability を作る．
\begin{mathpar}
\inferrule{ }{
  \DDPPat_{\Sighat}(q,S,\$,\tau,
    [\chi_{q_\kappa}\triangleright\tau],\varnothing,
    q+\langle1,0\rangle,S)
}\rn{DD-PP-HOLE}

\inferrule{ }{
  \DDPPat_{\Sighat}(q,S,\_,\tau,[],\varnothing,q,S)
}\rn{DD-PP-WILD}

\inferrule{ }{
  \DDPPat_{\Sighat}(q,S,\#\$x,\tau,[],x:\tau,q,S)
}\rn{DD-PP-VALUE}
\end{mathpar}

constructor と tuple は result を expected target に align してから子へ expected target 列を渡す．
$\mathsf{DDPPats}$ は hole 列を連結し，binder の disjointness を検査する．
\begin{mathpar}
\inferrule{
  \Inst_q(\Sighat_P(c))=(\overline\rho\Rightarrow_P\rho,q_0) \\
  \DDAlignTypes(S,\rho,\tau,S_1) \\
  \mathsf{DDPPats}_{\Sighat}(q_0,S_1,\overline{pp},\overline\rho,
    \overline\eta,\Delta,q',S')
}{
  \DDPPat_{\Sighat}(q,S,c\,\overline{pp},\tau,
    \overline\eta,\Delta,q',S')
}\rn{DD-PP-CON}

\inferrule{
  \mathsf{freshTargets}(n,q)=(\overline\alpha,q_0) \\
  \DDAlignTypes(S,(\overline\alpha),\tau,S_1) \\
  \mathsf{DDPPats}_{\Sighat}(q_0,S_1,\overline{pp},\overline\alpha,
    \overline\eta,\Delta,q',S')
}{
  \DDPPat_{\Sighat}(q,S,(\overline{pp}),\tau,
    \overline\eta,\Delta,q',S')
}\rn{DD-PP-TUPLE}
\end{mathpar}

$\DDDPat$ は arm の data pattern を clause target に対して check し，bindings を返す．
\begin{mathpar}
\inferrule{ }{
  \DDDPat_{\Sighat}(q,S,\$x,\tau,x:\tau,q,S)
}\rn{DD-DP-VAR}

\inferrule{ }{
  \DDDPat_{\Sighat}(q,S,\_,\tau,\varnothing,q,S)
}\rn{DD-DP-WILD}

\inferrule{
  \Inst_q(\Sighat_D(C))=(\overline\rho\to\rho,q_0) \\
  \DDAlignTypes(S,\rho,\tau,S_1) \\
  \mathsf{DDDPats}_{\Sighat}(q_0,S_1,\overline{dp},\overline\rho,
    \Delta,q',S')
}{
  \DDDPat_{\Sighat}(q,S,C\,\overline{dp},\tau,\Delta,q',S')
}\rn{DD-DP-CON}

\inferrule{
  \mathsf{freshTargets}(n,q)=(\overline\alpha,q_0) \\
  \DDAlignTypes(S,(\overline\alpha),\tau,S_1) \\
  \mathsf{DDDPats}_{\Sighat}(q_0,S_1,\overline{dp},\overline\alpha,
    \Delta,q',S')
}{
  \DDDPat_{\Sighat}(q,S,(\overline{dp}),\tau,\Delta,q',S')
}\rn{DD-DP-TUPLE}
\end{mathpar}

\subsection{Clause と matcher literal の型付け}

一つの clause は primitive pattern，next-matcher 列，arm 列の順に処理する．
$\decomposeME(M,n)$ は next expression を hole 数と同じ長さの expression 列へ展開する．
各 next matcher は対応する hole slot に対して $\DDCheck$ される．
\begin{mathpar}
\inferrule{
  \DDPPat_{\Sighat}(q,S,pp,\tau,\overline\eta,
    \Delta_{pp},q_1,S_1) \\
  \decomposeME(M,|\overline\eta|)=\mathsf{some}\;\overline M \\
  \DDChecks_{\Sighat}(q_1,S_1,\Gamma,\overline M,
    \overline{\MatcherSlot{\eta.\kappa}{\eta.\tau}},q_2,S_2) \\
  \DDArms_{\Sighat}(q_2,S_2,\Gamma,\Delta_{pp},\overline a,
    \tau,\ListTy{\prodty(\overline{\eta.\tau})},q',S')
}{
  \DDClause_{\Sighat}(q,S,\Gamma,
    pp\ \mathtt{as}\ M\ \mathtt{with}\ \overline a,
    \tau,\overline\eta,q',S')
}\rn{DD-CLAUSE}
\end{mathpar}
$\DDArms$ は data-pattern bindings と primitive-pattern bindings の衝突を拒否し，arm body を
$\ListTy{\prodty(\overline{\eta.\tau})}$ に対して check する．

matcher literal は全 clause で一つの fresh target を共有する．最後の substitution で hole dual を
一度だけ解決して evidence を再計算し，shape，catch-all，binder，arm exhaustiveness，coverage の
全 checker を通した場合だけ matcher type を返す．
\begin{mathpar}
\inferrule{
  \DDClauses_{\Sighat}(q+\langle0,1\rangle,S,\Gamma,\mathit{cls},
    \alpha_{q_\tau},
    \overline{\overline\eta},q',S') \\
  \mathsf{collectClauseEvidence}_{\Sighat}
    (\mathit{cls},\mathsf{terminalHoleCaps}(S',
      \overline{\overline\eta}))=\mathsf{some}\;\overline{ev} \\
  \mathsf{inferShape}_{\Sighat}(\overline{ev})=
    \mathsf{some}\;\kappa \\
  \mathsf{clauseCapsListCheck}_{\Sighat}
    (\kappa,\mathit{cls},\mathsf{terminalHoleCaps}
      (S',\overline{\overline\eta})) \\
  \CatchAllLast(\mathit{cls})\quad
  \mathsf{MatcherBindersCheck}(\mathit{cls}) \\
  \ArmExhaustive_{\Sighat_D}(\mathit{cls},S'\alpha_{q_\tau})\quad
  \CoverageOK_{\Sighat}(\mathit{cls},\kappa)
}{
  q;S;\Gamma\vdash
  \mathtt{matcher}\;\mathit{cls}\Rightarrow
  \Matcher{\kappa}{\alpha_{q_\tau}}\dashv q';S'
}\rn{DD-MATCHER}
\end{mathpar}

$\CatchAllLast$ は bare-hole catch-all を最後の clause に限定する．
$\ArmExhaustive$ は保守的で実行可能な checker であり，core では variable または wildcard arm
があれば成功する．$\CoverageOK$ は observable structured capability ごとに必要な constructor
clause があることを検査する．checker の失敗を warning に落とす mode は core に含めない．

matcher literal を body に持つ recursion は clause skeleton から domain／codomain placeholder を
決める専用規則を使う．
\begin{mathpar}
\inferrule{
  f\ne x\quad\DirectSelf(f,\mathtt{matcher}\;\mathit{cls}) \\
  \mathsf{fixMatcherPlaceholderSupply}_{\Sighat}
    (\mathit{cls},q)=\mathsf{some}\;(\tau_d,\tau_c,q_0) \\
  q_0;S;\Gamma,f:\mono{(\tau_d\to\tau_c)},x:\mono{\tau_d}
    \vdash\mathtt{matcher}\;\mathit{cls}\Rightarrow
      \rho\dashv q_1;S_1 \\
  \DDAlignTypes(S_1,\rho,\tau_c,S')
}{
  q;S;\Gamma\vdash
    \mathtt{fix}\ f\ x.(\mathtt{matcher}\;\mathit{cls})\Rightarrow
      \tau_d\to\tau_c\dashv q_1;S'
}\rn{DD-FIX-MATCHER}
\end{mathpar}
$\mathsf{NonMatcherBody}$ により通常の \textsc{DD-FIX} と排他的である．
